\documentclass[11pt]{article}
\usepackage[T1]{fontenc}
\usepackage{lmodern,microtype}
\usepackage{amsmath,amssymb,amsthm,mathtools}
\usepackage[letterpaper,margin=1in]{geometry}
\usepackage{booktabs,array,enumitem,longtable}
\usepackage{xcolor}
\usepackage[colorlinks=true,linkcolor=blue!55!black,citecolor=blue!55!black,urlcolor=blue!55!black]{hyperref}
\newcommand{\F}{F}
\newcommand{\K}{K}
\newcommand{\A}{\mathbb A}
\newcommand{\C}{\mathbb C}
\newcommand{\Z}{\mathbb Z}
\newcommand{\GL}{\operatorname{GL}}
\newcommand{\Irr}{\operatorname{Irr}}
\newcommand{\Nrd}{\operatorname{Nrd}}
\newcommand{\Nm}{\operatorname{N}}
\newcommand{\vol}{\operatorname{vol}}
\newcommand{\Gal}{\operatorname{Gal}}
\newcommand{\Aut}{\operatorname{Aut}}
\newcommand{\reg}{\mathrm{reg}}
\newcommand{\tH}{\widetilde H}
\newcommand{\tZ}{\widetilde Z}
\newcommand{\tT}{\widetilde T}
\newcommand{\Gb}{\overline G}
\newcommand{\GDb}{\overline{G_D}}
\newcommand{\Ptag}{\textnormal{\textbf{[P]}}}
\newcommand{\Etag}{\textnormal{\textbf{[E]}}}
\newcommand{\Ctag}{\textnormal{\textbf{[C]}}}
\newcommand{\Otag}{\textnormal{\textbf{[O]}}}
\theoremstyle{plain}
\newtheorem{theorem}{Theorem}[section]
\newtheorem{proposition}[theorem]{Proposition}
\newtheorem{lemma}[theorem]{Lemma}
\newtheorem{corollary}[theorem]{Corollary}
\theoremstyle{definition}
\newtheorem{definition}[theorem]{Definition}
\newtheorem{remark}[theorem]{Remark}
\newtheorem{example}[theorem]{Example}
\title{A metaplectic Jacquet--Langlands transfer for division algebras of prime degree, IV:\\
Wild local correspondence and global trace formula comparison\\[0.4em]\normalsize Draft 5}
\author{David Joyner \and ChatGPT (OpenAI) \and Claude (Anthropic)}
\date{November 7, 2026}
\emergencystretch=3em
\begin{document}
\maketitle
\begin{abstract}
We continue Parts I--III for an odd-degree covering of the multiplicative
inner form of $\GL_r$ with $r$ prime and $(r,n)=1$.  First we reconcile the
split-diagonal cocycle with the ambient-cover restriction and central
relabeling in Part I.  The regular good-element criterion for \emph{every}
maximal torus, including nonsplit tori and wild residue characteristic,
is already a theorem of Flicker--Kazhdan; we supply its proof with the
parameters used in this series.  We prove two refinements useful in global
comparison: an almost-everywhere local-to-global criterion for globally
good regular classes, and an exact local calculation of the kernel, index,
and Haar modulus of the $n$th-power map on an arbitrary finite \emph{\'etale}
torus modulo scalars.  The power-map calculation gives the split-place
root multiplicities, while a separate regularity example explains why a
regular root need not have a regular $n$th power.  The genuine character transform of Part III
is then placed alongside the corrected metaplectic character relation of
Zou and the older work of Flicker--Kazhdan.  A detailed wild rank-two example proves an integral, positive
Fourier decomposition of the discriminant ratio on every quadratic
maximal torus of the quaternion division algebra over
$\mathbb Q_3(\zeta_3)$, including both ramified and unramified
extensions; this is a torus-restriction result and does not prove
irreducible-character integrality for the full group. A general toral
Fourier criterion isolates the remaining nonabelian gluing question.
A normalization ledger and explicit Flicker--Kazhdan-compatible
split-place test-function conditions separate the previously established
covered-inner-form trace comparison from the additional proposed
comparison with ordinary $D^\times$.  A conditional geometric comparison is made explicit with its global
normalization factor, a necessary nongood-class support restriction,
and separate split-place matching conditions.  For $r=2$, an
unramified spherical-unit case follows from Flicker--Kazhdan in its
standard normalization, subject to transport to the specified ambient
cover. The wild irreducible-character integrality and the full global
comparison are not established here.
\end{abstract}
\tableofcontents

\section{Introduction and mathematical status}\label{sec:intro}
The problem concerns three different local representation theories.
Write $H_v=D_v^\times$, and let $\GDb_v$ denote its specified $n$-fold
central cover; the cover $\Gb_v$ of $\GL_r(K_v)$ is defined inside the
same ambient cover.  An additional \emph{ordinary} group $H_v$ enters
through the transform $L_{D,v}$ of Part III.  These are not interchangeable.
In particular, a comparison between $\GDb$ and $\Gb$ has already been
obtained in Parts~I--II (here cited as Draft~8 of each) under their
stated test-function and spectral hypotheses.  We cite Part~III
as Draft~9 throughout.  The new proposed comparison between $H$ and $\GDb$ is a
different assertion.

\paragraph{Global assumptions.}
Let $K$ be a totally imaginary number field with $\mu_n\subset K$,
$n\geq1$ odd.  Let $r\geq2$ be prime, $(r,n)=1$, let $L/K$ be cyclic
of degree $r$, and let $D$ be a central division $K$-algebra of
degree $r$ containing $L$.  Put $S=\{v:D_v\text{ is division}\}$.
Fix integers $r',c$ such that
\begin{equation}\label{eq:parameters}
 rr'\equiv1\pmod n,\qquad r-1+2c\equiv0\pmod n,
 \qquad c_v\equiv r'c\pmod n.
\end{equation}
Thus $r-1+2rc_v\equiv0\pmod n$.  In fact $c_v$ is independent
of $v$ with these choices; the subscript recalls Part I's notation.
Locally $F=K_v$ is a nonarchimedean field containing $\mu_n$;
$p$ denotes its residue characteristic.  We call $v$ \emph{wild}
if $p\mid n$.

\paragraph{Status convention.}
\textbf{[P]} denotes a proof included here (possibly invoking
\textbf{[E]} inputs established in cited publications or earlier parts);
\textbf{[O]} denotes a proof outlined but not fully established, and
\textbf{[C]} denotes a conjectural goal requiring new mathematical
input.  When a proved proposition has explicit unresolved hypotheses,
we mark its formal implication [P] and mark the outstanding hypotheses
[C] separately, in agreement with the status usage of Part III,
Draft~9.

\begin{center}\small
\begin{tabular}{@{}p{0.49\linewidth}p{0.37\linewidth}p{0.07\linewidth}@{}}
\toprule
Assertion & Precise basis & Status\\\midrule
Diagonal and nonsplit good-element criterion, including wild $p$ & [FK, \S3]; Part I, \S3 & [E]\\
Compatibility of the Part I cocycle with $\sigma_{c_v}$ & Part I, Prop. 2.3(b); direct calculation & [P]\\
Power-map kernel, index, and Haar modulus on $T_E(F)/F^\times$ & Section \ref{sec:localpowers} & [P]\\
Global good-regular class criterion & Section \ref{sec:globalgood}; Part III Hasse principle & [P]\\
Division-side unitary character transform and tame integral case & Part III, Thms. 4.3, 7.6 & [E]\\
Wild rank-two torus Fourier expansions & Worked example in \S\ref{subsec:examplewild} & [P]\\
Wild irreducible-character integrality for $\widetilde{D^\times}$ & Section \ref{sec:wild} & [C]\\
Covered $D$--covered $\GL$ global simple trace comparison & Parts I--II under their hypotheses & [E]\\
Covered $D$--ordinary $D$ global comparison & Conditional implication in \S\ref{sec:geometric}; full matching open & [P]/[C]\\
Removal of auxiliary support and packet hypotheses & Section \ref{sec:spectral} & [C]\\
\bottomrule
\end{tabular}
\end{center}

\subsection*{The running example}
The series' running example is $K=\mathbb{Q}(\zeta_3)$, $n=3$, $r=2$, $D=(-1,4-\zeta_3)_K$, ramified exactly at $S=\{\mathfrak p_3,\mathfrak p_7\}$, with $\mathfrak p_3=(1-\zeta_3)$ and $\mathfrak p_7=(3+\zeta_3)$ (Part~I, \S1, and \texttt{notes/mjl\_running\_example.tex}). Its wild place $\mathfrak p_3$ has $K_v=\mathbb{Q}_3(\zeta_3)$, and $D_{\mathfrak p_3}$ is the quaternion algebra of Example~\ref{ex:ranktwo-data}. Proposition~\ref{prop:all-three-tori} treats all three quadratic tori there. Theorem~\ref{thm:globalgood} describes the globally good classes, and Example~\ref{ex:ranktwoledger} is the global normalization ledger, in which $\Lambda=1$ forces $\lambda_\infty=1/3$. The smallest open case of wild integrality is whether $L_D\mathbf 1$ at $\mathfrak p_3$ is irreducible of dimension $27$.

\section{The exact cocycle convention and star map}\label{sec:cocycle}
This section fixes the one place where a parameter mismatch could
invalidate every subsequent constant.  The ambient construction is
a restriction of a cocycle over $L_v$, whereas [FK] uses the standard
$\GL_r(F)$ Kazhdan--Patterson convention.

\subsection{Restriction and relabeling}
Let $\sigma_F$ be the parameter-zero Kazhdan--Patterson cocycle with
restriction to the diagonal torus
\begin{equation}\label{eq:parameterzero}
\sigma_F(\operatorname{diag}(a_i),\operatorname{diag}(b_i))
=\prod_{i<j}(a_i,b_j)_{n,F}.
\end{equation}
The $\GL$ subgroup of the ambient cover in Part I has, \emph{as an
extension}, cocycle model
\begin{equation}\label{eq:before}
  \sigma_F(g,h)^r\,(\det g,\det h)_{n,F}^{c}
  \quad\text{before relabeling }\mu_n.
\end{equation}
This assertion is Part I, Proposition 2.3(b).  Applying the
kernel relabeling $\zeta\mapsto\zeta^{r'}$ yields the model
\begin{equation}\label{eq:after}
 \sigma_{c_v}(g,h)=\sigma_F(g,h)
 (\det g,\det h)_{n,F}^{c_v}.
\end{equation}
The two parameters $c$ and $c_v$ must be kept separate: $c$
occurs before relabeling and $c_v=r'c\pmod n$ afterward.
The isomorphism in \eqref{eq:before} is an isomorphism of extensions,
not a claim that arbitrary choices of cocycle sections give
literally equal functions on all of $\GL_r(F)^2$.

\begin{proposition}[\Ptag\ Diagonal commutator audit]\label{prop:diag}
For $t=\operatorname{diag}(t_1,\ldots,t_r)$ and
$u=\operatorname{diag}(u_1,\ldots,u_r)$ in $\GL_r(F)$,
\begin{equation}\label{eq:matrix}
 [\widetilde t,\widetilde u]
 =\prod_{i,j}(t_i,u_j)_{n,F}^{A_{ij}},\qquad
 A=-I_r+(1+2c_v)J_r.
\end{equation}
Here $J_r$ is the all-ones matrix.  Its row sums are
$r-1+2rc_v\equiv0\pmod n$, so scalar lifts commute with the lifted
split torus.  The radical on $(F^\times/F^{\times n})^r$ consists
precisely of diagonal classes.
\end{proposition}
\begin{proof}
The commutator is the ratio
$\sigma_{c_v}(t,u)/\sigma_{c_v}(u,t)$.
The Hilbert symbol satisfies $(a,b)_n=(b,a)_n^{-1}$; this is
\emph{antisymmetry}, valid for all $n$.  The determinant twist
contributes $2c_v$ to each entry of $A$, and the parameter-zero
factor contributes $1$ off the diagonal.  This gives \eqref{eq:matrix}.
The Hilbert pairing on $X=F^\times/F^{\times n}$ is perfect.
As $X$ is killed by $n$ and $(r,n)=1$, the decomposition
$X^r=\Delta X\oplus\{(x_i):\sum x_i=0\}$ holds by averaging with
$r^{-1}$ modulo $n$.  On these summands $A$ acts, respectively, by
$r-1+2rc_v=0$ and $-1$.  Therefore $\ker A=\Delta X$.
\end{proof}

\begin{remark}
The assumption $(r,n)=1$ also follows algebraically from
$n\mid r-1+2rc_v$, because any common divisor of $r$ and $n$ must
divide $1$.  We retain it because it is a global hypothesis of the
series.  The oddness of $n$ is used for the ordinary star-map
convention below, not for Hilbert-symbol antisymmetry.
\end{remark}

\subsection{Star and central characters}\label{subsec:star}
For an element $h$ of the ordinary group, choose any lift
$\widetilde h$ and define $h^*=\widetilde h^{\,n}$.  It is
independent of the lift because $\mu_n$ is central.  In the present
scalar-central cover, for $a\in F^\times$,
\begin{equation}\label{eq:starcentral}
 (ah)^*=a^*h^*,\qquad a^*=\widetilde a^{\,n}.
\end{equation}
If $h_1,h_2$ commute in an ordinary torus, their lifted commutator
is in $\mu_n$, and the class-two power identity
\begin{equation}\label{eq:starpower}
 (\widetilde h_1\widetilde h_2)^n
 =\widetilde h_1^{\,n}\widetilde h_2^{\,n}
 [\widetilde h_2,\widetilde h_1]^{n(n-1)/2}
 =\widetilde h_1^{\,n}\widetilde h_2^{\,n}
\end{equation}
uses $n$ odd.  Thus star is multiplicative on the lifted torus.
This agrees with the convention in [FK, \S4] and [Zou, \S3]:
Zou writes $x^*=u(x)s(x)^n$, where $u(x)=\pm1$ is the
correction defined in \cite[\S4]{FK} and \cite[\S3.1]{Zou24}.
For odd $n$ the extra sign is absent in the
Flicker--Kazhdan convention: the star is the $n$th power of a lift.
The correction $u(x)$ is introduced for the even-$n$ convention;
we do not infer that an arbitrary sign automatically lies in
$\mu_n$.  In the odd-degree setting, therefore, $u(x)=1$
by definition of the relevant convention (compare
\cite[\S4]{FK} and \cite[\S3.1]{Zou24}).  For a scalar $\zeta\in\mu_n(F)$,
$\zeta^*=1$ in the standard diagonal cocycle model: the
self-symbol $(\zeta,\zeta)_n=1$ for odd $n$, so
$s(\zeta I_r)^n=s(I_r)=1$.  This last statement is also
invariant under the extension isomorphisms in Part I.
For representation characters we first fix a genuine unitary
character $\widetilde\theta$ of the central preimage $\tZ$ and
\emph{then} put $\theta(a)=\widetilde\theta(a^*)$.
The reverse implication is not unique: $\theta$ alone does not
specify $\widetilde\theta$.  For $\zeta\in\mu_n(F)$ we have
$\zeta^*=1$, hence $\theta(\zeta)=1$.  Globally the rational
splitting gives $a^*=\iota(a^n)$ for $a\in K^\times$;
if $\widetilde\theta$ is automorphic, it is trivial on the rational
splitting, so $\theta(a)=1$ on $K^\times$.
Test functions use the inverse central characters
$\theta^{-1}$ and $\widetilde\theta^{-1}$.

\section{Good elements for every maximal torus}\label{sec:good}
A regular maximal torus in $\GL_r(F)$ may be nonsplit even at a
place where $D_v\simeq M_r(F)$.  This is the situation not treated by
the diagonal matrix alone.  The following calculation is not new:
it is the proof of [FK, \S3, Proposition] in the notation of this
series, and also underlies Part I, Proposition 3.5.

\begin{theorem}[\Etag\ Flicker--Kazhdan good-element theorem]\label{thm:FKgood}
Let $E=\prod_{i=1}^kE_i$ be a finite \emph{\'etale} $F$-algebra
of dimension $r$, embedded as a regular maximal torus
$T_E(F)=E^\times$ in $\GL_r(F)$.  Under \eqref{eq:parameters},
for any $t\in T_E(F)$ and any lift $\widetilde t$,
\begin{equation}\label{eq:FKtorus}
 [\widetilde t,\widetilde u]
 =(\Nm_{E/F}t,\Nm_{E/F}u)_{n,F}^{1+2c_v}
   \prod_{i=1}^k(t_i,u_i)_{n,E_i}^{-1}.
\end{equation}
The radical of this commutator pairing is
\begin{equation}\label{eq:everygood}
 \{t\in T_E(F):[\widetilde t,\widetilde u]=1\ \forall u\in T_E(F)\}
 =F^\times T_E(F)^n.
\end{equation}
For regular $t$, this is exactly the criterion that $\widetilde t$
is good.  Both identities are valid when $p\mid n$.
\end{theorem}
\begin{proof}[Proof from the cited formula, included for clarity]
The commutator formula is [FK, \S3], using [KP, Prop.~0.1.5];
its restriction to a single field is exactly Part I,
Proposition 3.1 after relabeling \cite{JoynerI}.  Write $N=\Nm_{E/F}$.
Norm compatibility of local Hilbert symbols gives
\[
 (Nt,Nu)_{n,F}
 =\prod_i(Nt,u_i)_{n,E_i}.
\]
Thus \eqref{eq:FKtorus} equals
\[
 \prod_i\bigl((Nt)^{1+2c_v}t_i^{-1},u_i\bigr)_{n,E_i}.
\]
The nondegeneracy of the $n$th Hilbert symbol over each $E_i$
(including wildly ramified $E_i$) implies that the commutator is
trivial for all $u$ exactly when
\begin{equation}\label{eq:pointwisegood}
 t_i\equiv(Nt)^{1+2c_v}\pmod{E_i^{\times n}}
 \quad (1\leq i\leq k).
\end{equation}
This directly says $t\in F^\times E^{\times n}$.  Conversely,
if $t=ay^n$ for $a\in F^\times$, then
$Nt=a^r(Ny)^n$; therefore
$(Nt)^{1+2c_v}\equiv a^{r(1+2c_v)}\equiv a\equiv t_i$
modulo $E_i^{\times n}$ by scalar centrality.  This proves the
criterion.  No tameness assumption has been introduced.
\end{proof}

\begin{example}[The diagonal rank-two check]\label{ex:diag}
Take $r=2$, $n=3$ and $F=\mathbb Q_3(\zeta_3)$.
In the notation of Part I, Example 3.7, choose $c=1$, so
$c_v\equiv2\pmod3$.  Formula \eqref{eq:matrix} becomes
\[
 A\equiv\begin{pmatrix}1&-1\\-1&1\end{pmatrix}\pmod3,
 \qquad
 [\widetilde t,\widetilde u]
 =(t_1/t_2,u_1/u_2)_{3,F}.
\]
Thus $\operatorname{diag}(t_1,t_2)$, when regular, is good exactly
when $t_1/t_2\in F^{\times3}$.  If a quadratic field $E/F$ is
embedded as a nonsplit maximal torus, the same theorem says that
goodness is equivalent to $t\in F^\times E^{\times3}$; the two
tests are instances of the same formula, not separate conjectures.
\end{example}

\begin{remark}[Why goodness is not enough for a regular-root formula]\label{rem:power-singular}
The power of a regular element in a \emph{split} torus need not be
regular.  For $n=3$ and $F=\mathbb Q_3(\zeta_3)$, the element
$\operatorname{diag}(1,\zeta_3)$ is regular but not good:
$\zeta_3\notin F^{\times3}$, since $F$ does not contain $\zeta_9$.
Its cube is the identity.  A character or orbital-integral formula with
$\Delta(t^n)$ in the denominator must restrict to $t^n$ regular
(as in Zou's equation (3.2)), or use a separately proved germ extension.
On a degree-$r$ \emph{field} torus, by contrast, if $t$ is regular
and $t^n\in F$, then Kummer theory gives
$[F(t):F]\mid n$; since $[F(t):F]=r$ and $(r,n)=1$, this is
impossible.  Because $r$ is prime, $F(t^n)=F(t)$ and $t^n$ is regular.
\end{remark}

\section{Local power-map calculation for \texorpdfstring{\'etale}{etale} tori}\label{sec:localpowers}
The key local distinction is between a field centralizer (compact
modulo scalars) and a product of fields (not compact modulo scalars).
Both occur in the split \emph{group}.  The next proposition gives
an exact count of the roots and the size of the image at wild and
tame places, including the factors that a global trace identity
must account for.

\begin{lemma}[\Ptag\ Kummer injection for a product algebra]\label{lem:Kumprod}
Let $E=\prod_{i=1}^kE_i$ be finite \emph{\'etale} of dimension $r$
over $F$, with $\mu_n\subset F$ and $(r,n)=1$.  Then
\begin{equation}\label{eq:Kummerinj}
 F^\times/F^{\times n}\longrightarrow
 E^\times/E^{\times n}
\end{equation}
is injective.
\end{lemma}
\begin{proof}
Suppose $a\in F^\times$ is an $n$th power in every $E_i$.
The Kummer extension $F(a^{1/n})/F$, since $\mu_n\subset F$,
is cyclic of degree $d\mid n$ and embeds in every $E_i$.
Therefore $d\mid[E_i:F]$ for every $i$, so also
$d\mid\sum_i[E_i:F]=r$.  Coprimality gives $d=1$ and
$a\in F^{\times n}$.
\end{proof}

\begin{proposition}[\Ptag\ Kernel, index, and Haar modulus]\label{prop:power-etale}
Under the hypotheses of Lemma \ref{lem:Kumprod}, put
$A_E=E^\times/F^\times$ and $P_n([t])=[t^n]$.  Then
\begin{align}
 \ker(P_n)&\simeq \mu_n(E)/\mu_n(F)
      \simeq \mu_n^{\,k}/\mu_n,
 &|\ker(P_n)|&=n^{k-1},\label{eq:kernel}\\
 [A_E:P_n(A_E)]&=n^{2(k-1)}|n|_F^{-(r-1)}.
 \label{eq:index}
\end{align}
In particular, $P_n(A_E)$ is open and closed.  With \emph{any}
Haar measure $da$ on $A_E$ and its restriction to $P_n(A_E)$,
for compactly supported integrable $q$ on that image,
\begin{equation}\label{eq:haarJac}
 \int_{A_E}q(a^n)\,da
 =n^{k-1}|n|_F^{-(r-1)}
   \int_{P_n(A_E)}q(b)\,db.
\end{equation}
\end{proposition}
\begin{proof}
If $t^n=a\in F^\times$, then $a$ is an $n$th power in every
$E_i$; by Lemma \ref{lem:Kumprod}, write $a=b^n$ with $b\in F^\times$.
Thus each $t_i/b\in\mu_n(E_i)=\mu_n(F)$, giving
\eqref{eq:kernel}, including its canonical description.

For a local field $M$ containing $\mu_n$ the Kummer quotient has
\begin{equation}\label{eq:localKumIndex}
 |M^\times/M^{\times n}|=n^2|n|_M^{-1}.
\end{equation}
Indeed the discrete valuation contributes $n$; the finite
torsion unit factor containing $\mu_n$ contributes $n$; and
the torsion-free $p$-adic unit factor contributes
$|n|_M^{-1}$ (multiplication by $n$ on its $\Z_p$-lattice).  Alternatively, \eqref{eq:localKumIndex}
follows from local duality.
Since $F^\times/F^{\times n}$ injects into
$E^\times/E^{\times n}$, the group $A_E/A_E^n$ is its cokernel:
\[
 |A_E/A_E^n|
 =\frac{\prod_i n^2|n|_{E_i}^{-1}}
             {n^2|n|_F^{-1}}
 =n^{2(k-1)}|n|_F^{-(r-1)},
\]
where $|n|_{E_i}=|n|_F^{[E_i:F]}$ for normalized absolute values
on local fields.  This proves \eqref{eq:index}.

For the Haar statement, write $K_0=\ker(P_n)$, of order
$n^{k-1}$.  Choose a sufficiently small compact open subgroup
$U\subset A_E$ consisting of principal-unit classes and having
$U\cap K_0=\{1\}$.  The $p$-adic logarithm identifies $U$ with
a free $\mathbb Z_p$-module of rank
$(r-1)[F:\mathbb Q_p]$: the factors $E_i^\times$ have combined
principal-unit rank $r[F:\mathbb Q_p]$, and quotienting by the
scalar principal units subtracts $[F:\mathbb Q_p]$.
Multiplication by $n$ maps $U$ injectively into itself, with
\[
 [U:U^n]=p^{v_p(n)(r-1)[F:\mathbb Q_p]}
        =|n|_F^{-(r-1)}.
\]
Since $P_n^{-1}(U^n)=K_0U$ and $K_0\cap U=\{1\}$, its
pushforward of Haar measure gives
\[
 \frac{\operatorname{vol}(P_n^{-1}(U^n))}
      {\operatorname{vol}(U^n)}
   =\frac{|K_0|\operatorname{vol}(U)}
          {\operatorname{vol}(U)/[U:U^n]}
   =n^{k-1}|n|_F^{-(r-1)}.
\]
The pushforward is translation-invariant on the open image
$P_n(A_E)$, so this ratio is its constant density on that image,
proving \eqref{eq:haarJac}.  Notice that this argument requires
neither a global quotient-index formula nor a choice of scale for
Haar measure.
\end{proof}

\begin{remark}[A measure-theoretic warning]\label{rem:haar}
For a general open endomorphism $\varphi:A\to A$ of a locally
compact abelian group with finite kernel, its pushforward Haar density
need not equal $[A:\varphi A]/|\ker\varphi|$.  For example, on
$A=\mathbb Z$ with counting measure, $a\mapsto na$ has density
$1$ on $n\mathbb Z$, not $n$.  The correct local computation here is
$|\ker P_n|[U:U^n]$, where $U$ is the compact open subgroup in
the proof.  For this specific $A_E$ it happens to simplify to
$[A_E:P_nA_E]/|\ker P_n|$; that numerical equality is not a
general change-of-variables theorem.
\end{remark}

\begin{example}[Three local torus types]
If $E/F$ is a field of degree $r$, then $k=1$ and $P_n$ is
injective, with index and Haar modulus $|n|_F^{-(r-1)}$.
If $E=F^r$ then $k=r$ and the kernel has order $n^{r-1}$.
For $r=2,n=3$, $F=\mathbb Q_3(\zeta_3)$ and $|3|_F=1/9$,
\[
\begin{array}{c|ccc}
\text{type}&|\ker P_3|&[A_E:A_E^3]&\text{factor in \eqref{eq:haarJac}}\\\hline
E/F\text{ quadratic}&1&9&9\\
E=F\times F&3&81&27
\end{array}
\]
These integers quantify why one cannot use the division-torus
substitution formula unchanged on the full split group.
\end{example}

\section{A global classification of good regular classes}\label{sec:globalgood}
The local theorem now applies to the actual torus at \emph{every}
place.  Combining it with the Hasse principle of Part III removes
the local-to-global ambiguity for power classes of a degree-$r$
field, though not the independent measure and spectral questions.

\begin{proposition}[\Etag\ Almost-everywhere powers modulo scalars]\label{prop:Hasse}
Let $E/K$ be a field extension of degree $r$ with $(r,n)=1$ and
$\mu_n\subset K$.  For every finite set $\Sigma$ of places,
\begin{equation}\label{eq:Hasse}
 E^\times/(K^\times E^{\times n})\longrightarrow
 \prod_{v\notin\Sigma}
 \frac{(E\otimes_K K_v)^\times}
 {K_v^\times(E\otimes_KK_v)^{\times n}}
\end{equation}
is injective.  This is Part III, Proposition 11.3 \cite{JoynerIII}.
\end{proposition}
\begin{proof}[Recalled proof]
Let $t$ be locally scalar times an $n$th power outside $\Sigma$.
Choose $u,k\in\Z$ with $ru+nk=1$ and put
$a=N_{E/K}(t)^u\in K^\times$.  If $t=a_v y_v^n$ in the local
\'etale algebra, taking norms gives the identity
\[
 t/a=\bigl(a_v^k y_v/N_{E_v/K_v}(y_v)^u\bigr)^n.
\]
Thus $t/a$ is an $n$th power in every completion of $E$ over all
$v\notin\Sigma$.  If its class in $E^\times/E^{\times n}$ were
nontrivial, the nontrivial Kummer extension
$E((t/a)^{1/n})/E$ would, by Chebotarev, admit infinitely many
unramified places at which $t/a$ is not an $n$th power.
This contradicts the hypothesis.  Hence $t\in K^\times E^{\times n}$.
\end{proof}

\begin{theorem}[\Ptag\ Global good-regular criterion]\label{thm:globalgood}
Assume the covers and rational splittings of Part I.  Let
$G$ denote $D^\times$ or $\GL_r$, and let $\gamma\in G(K)$ be
regular elliptic, so $E=K[\gamma]$ is a field of degree $r$.
The following are equivalent:
\begin{enumerate}[label=(\alph*)]
\item the rational lift $\iota(\gamma)$ is good at every place;
\item it is good at every place outside some finite set $\Sigma$;
\item $\gamma\in K^\times E^{\times n}$.
\end{enumerate}
The map $[s]\mapsto[s^n]$ is a bijection from the noncentral
classes in $E^\times/K^\times$ onto the globally good regular
classes belonging to $E$.  Furthermore
\begin{equation}\label{eq:Wstabilizer}
 W_{s^n}=W_s,\qquad
 W_s=\{\sigma\in\Aut(E/K):\sigma(s)\in K^\times s\}.
\end{equation}
\end{theorem}
\begin{proof}
At each finite $v$, the ordinary centralizer is
$T_v=(E\otimes_K K_v)^\times$.  If $D_v$ splits, apply
Theorem \ref{thm:FKgood}, transported through Part I, Proposition~2.6 (the split-place
isomorphism of the specified covers).  If $D_v$ is division, use Part I,
Proposition 3.5(a), whose commutator is the same by Part I,
Proposition 3.1.  At complex places the $n$th-power map on
$T_v$ is surjective.  Thus at every place the local goodness
condition is precisely
$\gamma\in K_v^\times T_v^n$.
Implications (c)$\Rightarrow$(a)$\Rightarrow$(b) follow,
while (b)$\Rightarrow$(c) is Proposition \ref{prop:Hasse}.

The power map $E^\times/K^\times\to E^\times/K^\times$
is injective: if $s^n\in K^\times$, Kummer theory shows
$[K(s):K]\mid n$, while $[K(s):K]\mid [E:K]=r$.
Hence $s\in K^\times$.  Because $r$ is prime, every noncentral
$s\in E$ generates $E$, and the same is true of $s^n$.
This proves the bijection.  Finally
$\sigma(s^n)/s^n\in K^\times$ if and only if
$(\sigma(s)/s)^n\in K^\times$, which is equivalent by injectivity
to $\sigma(s)/s\in K^\times$.  Equation \eqref{eq:Wstabilizer}
follows.
\end{proof}

\begin{remark}[Precise consequence for the geometric side]
Part I, Proposition 6.1, uses the rational regular classes which
are good everywhere, with coefficient
\begin{equation}\label{eq:globalvol}
 c(\gamma)=\frac{\vol\bigl(\A_K^\times E^\times\backslash
 \A_E^\times;dt\bigr)}{|W_\gamma|}.
\end{equation}
Theorem \ref{thm:globalgood} parametrizes its noncentral
summation set by global $[s^n]$ and preserves $W_\gamma$.
It does \emph{not} replace the local root sums in Zou's
character formula: the global root class may be unique even when
a localized split torus has $n^{k_v-1}$ roots modulo scalars.
\end{remark}

\section{The remaining wild local character problem}\label{sec:wild}
We distinguish character transfer from class-function transfer.
The latter is a theorem of Part III; the former requires an
\emph{integral} statement not implied by unitarity.

\subsection{Division-side normalization and discriminants}
Let $F$ be nonarchimedean, $D/F$ a division algebra of prime
degree $r$, $H=D^\times$, and $\tH$ its specified cover.
Set $Z=F^\times$ and $\tZ=p^{-1}(Z)$, equipped with compatible
Haar measures on $H/Z\simeq\tH/\tZ$.  The Part III transform is
\begin{equation}\label{eq:LD}
 (L_D\Phi)(z\,t^*)
 =b_D\,\widetilde\theta(z)
  \frac{\Delta(t)}{\Delta(t^n)}\Phi(t),\qquad
 b_D=|n|_F^{-(r-1)/2},
\end{equation}
for regular $t$, $z\in\tZ$; it vanishes at nongood regular lifts.
Here
\begin{equation}\label{eq:Delta}
 \Delta(t)=\left|\det\bigl(1-\operatorname{Ad}(t)
 \mid\mathfrak h/\mathfrak t\bigr)\right|_F^{1/2}
 =\prod_{i<j}\left|
 \frac{(\lambda_i-\lambda_j)^2}{\lambda_i\lambda_j}
 \right|_F^{1/2}.
\end{equation}
This equals Part I's $|D(t)|^{1/2}|\Nrd(t)|_F^{-(r-1)/2}$.
Since the power map on a field torus modulo scalars has index
$|n|_F^{-(r-1)}$, Weyl integration gives $b_D^2$
equal to this index.  The factor $b_D$ is independent of
\emph{common} rescalings of the two identified quotient Haar
measures.  The ratio of discriminants extends at central elements
to $|n|_F^{-r(r-1)/2}$ (Part III, Lemma 3.3), hence
\begin{equation}\label{eq:cD}
 (L_D\Theta_\rho)(1)=c_D\dim\rho,\qquad
 c_D=b_D|n|_F^{-r(r-1)/2}
 =|n|_F^{-(r^2-1)/2}.
\end{equation}
The identity factor is pointwise and does not depend on Haar
measures.  Part III, Theorem 9.1, separately gives the adjoint
orbital-integral/test-function transfer.  Test functions transform
by $\widetilde\theta^{-1}$, not by $\widetilde\theta$.

\subsection{The integral matrix}\label{subsec:integral}
Choose $\widetilde\theta$ first and put
$\theta(a)=\widetilde\theta(a^*)$.
Part III, Theorem 4.3 and Corollary 5.1, show that $L_D$ is a
unitary \emph{isomorphism} of invariant $L^2$ spaces and that
\begin{equation}\label{eq:matrixchars}
 L_D\Theta_\rho=\sum_\tau a_{\rho\tau}\Theta_\tau,
 \qquad a_{\rho\tau}=\langle L_D\Theta_\rho,
 \Theta_\tau\rangle,
\end{equation}
with each row and column finite.  The characters on either side
are orthonormal bases in their fixed-central-character spaces.

\begin{proposition}[\Etag\ Integral-lattice criterion, Part III]\label{prop:int}
The transform $L_D$ permutes the irreducible characters with
positive signs if and only if every $a_{\rho\tau}$ in
\eqref{eq:matrixchars} is an integer.
\end{proposition}
\begin{proof}
An integral unit vector in $\ell^2$ is a signed basis vector;
unitarity gives injectivity and surjectivity because the inverse
is the adjoint of a unitary map.  The value at identity in
\eqref{eq:cD} is positive and fixes the sign.
This is Part III, Corollary 5.2, with its converse.
\end{proof}

\begin{remark}[The actual obstruction]
For wild $v$, Part III does not prove
$a_{\rho\tau}\in\Z$.  The tame Adams argument works with a
finite quotient $Q$ and a central subgroup $C$, under the
hypothesis $\gcd(n,|Q/C|)=1$ (Part III, Lemma 7.5).
At a wild place that hypothesis need not hold.  For the
simplest illustration take $Q=C_p$, $C=1$, $n=p$: then the
Adams operation sends every linear character to the trivial
character.  This does not prove a nonintegral coefficient for
$L_D$; it only shows why that particular proof method fails.
\end{remark}

\subsection{A fully worked wild example: all quadratic field tori}
\label{subsec:examplewild}
The following calculation is more informative than the dimension
obstruction alone.  It verifies a \emph{necessary representation-ring
integrality test on every maximal torus}, including the two ramified
quadratic field types; the full-group integrality statement remains open.

\begin{example}[The field, algebra, and cover]\label{ex:ranktwo-data}
Take
\[
 F=\mathbb Q_3(\zeta_3),\qquad \varpi=\zeta_3-1,\qquad
 3=-\zeta_3^{-1}\varpi^2,\qquad q_F=3.
\]
Thus $v_F(3)=2$ and $|3|_F=3^{-2}=1/9$.  The compatible
covering parameters in [JC] and Part~I are respectively
$c=1$ and $c_v=2\pmod3$.  Let $D/F$ be the unique quaternion
division algebra.  With $E_0=F(i)$, $i^2=-1$, its unramified
quadratic maximal subfield, we may write
\[
 D=E_0\oplus E_0j,\qquad j^2=\varpi,\qquad
 jx=\sigma(x)j\quad(x\in E_0),
\]
where $\sigma$ is the nontrivial automorphism of $E_0/F$.
Part~I, Proposition~3.5, gives, for every quadratic maximal
subfield $E/F$, the covered-torus commutator
\begin{equation}\label{eq:ranktwo-comm}
 [x,y]_{\widetilde{E^\times}}=
 (N_{E/F}x,N_{E/F}y)_{F,3}^{\,2}
 (x,y)_{E,3}^{-1}.
\end{equation}
The center $C_E$ of $\widetilde{E^\times}$ is the preimage of
$F^\times E^{\times3}$, and
\[
 [E^\times:F^\times E^{\times3}]=9,
 \qquad b_D=3,\qquad c_D=27.
\]
By the Heisenberg-type Stone--von Neumann theorem
([JC], Proposition~2.4; Part~I, Proposition~3.5), every genuine
irreducible representation of $\widetilde{E^\times}$ has degree $3$.
This is \emph{not} a degree calculation for representations of
$\widetilde{D^\times}$.
\end{example}

\begin{lemma}[A general toral Fourier test]\label{lem:toral-Fourier}
Let $E/F$ be any maximal field in $D$, with $r$ prime and
$(r,n)=1$, and put $A_E=E^\times/F^\times$.  The function
\[
 R_E(t)=\frac{\Delta(t)}{\Delta(t^n)}
\]
extends continuously at scalars, is a locally constant function
on the compact abelian group $A_E$, and has a finite Fourier expansion
\begin{equation}\label{eq:R-Fourier}
 R_E(t)=\sum_{\chi\in\widehat{A_E}}m_\chi\chi(t),
 \qquad
 m_\chi=\int_{A_E}R_E(t)\overline{\chi(t)}\,dt.
\end{equation}
Here the Haar measure of $A_E$ has mass $1$.
Choose a compatible genuine central character
$\widetilde\theta$ of $\widetilde Z$ and set
$\theta(a)=\widetilde\theta(a^*)$.  For any smooth character
$\alpha:E^\times\to\mathbb C^\times$ with
$\alpha|_{F^\times}=\theta$, and for each $\chi\in\widehat{A_E}$,
let $\eta_{\alpha\chi}$ be the unique genuine irreducible
representation of $\widetilde{E^\times}$ whose central character
satisfies
\[
 \zeta_{\alpha\chi}(z\,t^*)=
 \widetilde\theta(z)\alpha(t)\chi(t)
 \qquad(z\in\widetilde Z,\ t\in E^\times).
\]
Define a torus-level class function $\mathcal L_E\alpha$ to equal
$b_D\widetilde\theta(z)R_E(t)\alpha(t)$ on $z t^*\in C_E$
and to vanish outside $C_E$.  Then, as a genuine class function
on the \emph{covered torus},
\begin{equation}\label{eq:toral-branch}
 \mathcal L_E\alpha
 =\sum_\chi m_\chi\,\Theta_{\eta_{\alpha\chi}}.
\end{equation}
For an ordinary character $\Theta_\rho$ of $D^\times$,
its actual $L_D$ transform restricts to the sum of these
$\mathcal L_E\alpha$ with the multiplicities in
$\Theta_\rho|_{E^\times}$.
Consequently $m_\chi\in\mathbb Z_{\geq0}$ for every $\chi$
is sufficient for the restriction of every transferred ordinary
irreducible character to this covered torus to be a genuine
representation character.  It is not sufficient for the full
$\widetilde{D^\times}$ character to be a representation character.
\end{lemma}
\begin{proof}
The extension of $R_E$ at scalars follows from Part~III,
Lemma~3.3.  Since $E/F$ is a degree-$r$ field extension with
$(r,n)=1$, the $n$th power of a noncentral $t$ stays noncentral.
The compactness of $A_E$ and local constancy imply that $R_E$
factors through a finite quotient, giving \eqref{eq:R-Fourier}.

By Part~I the center of the covered torus is
$C_E=\widetilde Z\cdot\{t^*:t\in E^\times\}$.  The displayed
formula for $\zeta_{\alpha\chi}$ is well-defined: the
power map on $A_E$ is injective, and on scalar overlaps the
compatibility $\alpha(a)=\widetilde\theta(a^*)$ is precisely
what is required.  Such genuine central characters extend over
$C_E$, and Stone--von Neumann gives a unique $\eta_{\alpha\chi}$.
Its character vanishes off $C_E$, while on $C_E$ it equals
$b_D\zeta_{\alpha\chi}$ because its degree is $b_D$.
The defining formula \eqref{eq:LD} for $L_D$ has the identical
vanishing and equals $b_D\widetilde\theta(z)R_E(t)\alpha(t)$
on $z t^*$.  Comparing with \eqref{eq:R-Fourier} proves
\eqref{eq:toral-branch}.  Finally, restricting a finite-dimensional
ordinary representation to $E^\times$ decomposes it into
one-dimensional characters, all restricting to $\theta$ on
$F^\times$; summing the displayed formula proves the last claim.
\end{proof}

\subsubsection{The unramified quadratic torus}
Put $E_0=F(i)$ as above.  Since $X^2+1$ is irreducible over
$\mathbb F_3$, $E_0/F$ is unramified and its residue field is
$\mathbb F_9$.  For $t=a+bi$ with $a,b\in F$ the reduced norm
and trace give the elementary identity
\begin{equation}\label{eq:quadratic-R}
 R_{E_0}(a+bi)=
 \frac{|a^2+b^2|_F}{|3a^2-b^2|_F}.
\end{equation}
Multiplying $t$ by an element of $F^\times$, suppose
$a,b\in\mathcal O_F$ and at least one is a unit.
Then $a^2+b^2$ is always a unit.  If $b$ is a unit, so is
$3a^2-b^2$.  If $b=\varpi u$, then
\[
 3a^2-b^2=3(a^2+\zeta_3u^2),
\]
and $a$ is a unit; the reduction of $a^2+\zeta_3u^2$ modulo
$\varpi$ is either $1$ or $2$, never $0$.  Hence
\begin{equation}\label{eq:unram-values}
 R_{E_0}(a+bi)=
 \begin{cases}
 9,&b\in\varpi\mathcal O_F,\\
 1,&b\in\mathcal O_F^\times.
 \end{cases}
\end{equation}
The finite residue quotient is
\[
 B_0=k_{E_0}^\times/k_F^\times
 \simeq C_4.
\]
The condition $b\in\varpi\mathcal O_F$ says exactly that the
image of $t$ in $B_0$ is the identity.  Denote its four characters
by $\chi_0=1,\chi_1,\chi_2,\chi_3$.  The regular character of
$C_4$ equals $4$ at the identity and $0$ elsewhere. Thus
\begin{equation}\label{eq:unram-fourier}
 \boxed{\quad
 R_{E_0}=1+2\,\mathrm{reg}_{C_4}
 =3\chi_0+2\chi_1+2\chi_2+2\chi_3.
 \quad}
\end{equation}
This is a positive integral Fourier decomposition of degree
$3+2+2+2=9$; it is stronger than merely checking that
$R_{E_0}$ takes integer values.

\subsubsection{Both ramified quadratic tori}
Every ramified quadratic extension of $F$ can be written
$E_d=F(j_d)$, $j_d^2=d=u\varpi$ with $u\in\mathcal O_F^\times$.
There are two isomorphism classes, distinguished by the square
class of $u$.  For $t=a+bj_d$,
\begin{equation}\label{eq:ramified-R}
 R_{E_d}(a+bj_d)
 =\frac{|a^2-db^2|_F}{|3a^2+db^2|_F}.
\end{equation}
Again normalize $a,b$ to be integral and not both divisible
by $\varpi$.  The terms $3a^2$ and $db^2$ have different
valuations whenever both $a$ and $b$ are units.  A direct
valuation calculation therefore gives
\begin{equation}\label{eq:ramified-values}
 R_{E_d}(a+bj_d)=
 \begin{cases}
 9,& a\in\mathcal O_F^\times,
       \quad b\in\varpi\mathcal O_F,\\
 3,& a,b\in\mathcal O_F^\times,\\
 1,& a\in\varpi\mathcal O_F,
       \quad b\in\mathcal O_F^\times.
 \end{cases}
\end{equation}
For example, in the original cyclic algebra $d=\varpi$,
\[
 \begin{array}{c|ccc}
 t&j&1+j&1+\varpi j\\\hline
 R(t)&1&3&9\\
 (L_D1)(t^*)&3&9&27
 \end{array}
\]
for the central-character normalization of Example~\ref{ex:ranktwo-data}.
The three values illustrate why the unramified torus alone is not
an adequate test.

There is a natural finite quotient
\[
 B_d=E_d^\times/(F^\times U_1)\simeq C_2\times C_3,
\]
where $U_1\subset\mathcal O_{E_d}^\times$ consists of
$a+bj_d$ with $a\in\mathcal O_F^\times$ and
$b\in\varpi\mathcal O_F$.  Modulo $F^\times$ its
members are represented by $1+cj_d$ with
$c\in\varpi\mathcal O_F$.  The $C_2$ coordinate
records the parity of the $E_d$-valuation, and on even elements
the $C_3$ coordinate is $\bar c\in\mathbb F_3$ for
$1+cj_d$.  The multiplication law
\[
 (1+cj_d)(1+c'j_d)
 =(1+dcc')\left(1+\frac{c+c'}{1+dcc'}j_d\right)
\]
shows that $\bar c$ is additive, and $j_d^2=d\in F^\times$
provides the order-two factor.

Write $\epsilon$ for the nontrivial character of $C_2$
and $\xi$ for a character of order $3$ of $C_3$.
The identity \eqref{eq:ramified-values} says that $R_{E_d}$
has values $9$ at the identity of $B_d$, $3$ at its two
nonidentity elements in $C_3$, and $1$ on its three odd elements.
Consequently
\begin{equation}\label{eq:ram-fourier}
 \boxed{\quad
 \begin{aligned}
 R_{E_d}
 &=1+\mathrm{reg}_{C_2}+\mathrm{reg}_{C_2\times C_3}\\
 &=3\cdot1+2\epsilon
   +\xi+\xi^2+\epsilon\xi+\epsilon\xi^2.
 \end{aligned}\quad}
\end{equation}
Here the regular representations are inflated from the indicated
quotients.  All coefficients are nonnegative integers and sum
to $3+2+1+1+1+1=9$.  This applies to \emph{both} ramified
quadratic isomorphism classes, independently of $u$.

\begin{proposition}[\Ptag\ All quadratic torus restrictions pass
an integral-character test]\label{prop:all-three-tori}
In the wild example $F=\mathbb Q_3(\zeta_3)$, $r=2$, $n=3$,
let $D/F$ be the quaternion division algebra.  The function
$R_E$ is the character of a nine-dimensional permutation
representation of a finite quotient of $E^\times/F^\times$,
for every quadratic maximal subfield $E\subset D$.
More precisely, for unramified $E$ it is
$1\oplus 2\,\mathbb C[C_4]$, and for either ramified $E$
it is $1\oplus\mathbb C[C_2]\oplus\mathbb C[C_2\times C_3]$.
For a suitable genuine central character inducing the trivial
ordinary central character, the restriction of $L_D\mathbf1$
to $\widetilde{E^\times}$ is therefore an \emph{actual genuine
representation character} of total dimension $27$, namely
\[
 \begin{array}{c|c}
 E/F&\text{multiplicities of three-dimensional genuine torus irreducibles}\\\hline
 \text{unramified}&(3,2,2,2)\\
 \text{ramified, either type}&(3,2,1,1,1,1).
 \end{array}
\]
These multiplicities are invariant under the nontrivial
$F$-automorphism of each $E$.
\emph{No irreducibility or even integrality for the full
$\widetilde{D^\times}$ is asserted.}
\end{proposition}
\begin{proof}
A $p$-adic field of odd residue characteristic has four
square classes; hence $F$ has exactly three quadratic field
extensions, one unramified and two ramified.  Every quadratic
field extension embeds in the quaternion division algebra.
Equations \eqref{eq:unram-fourier} and \eqref{eq:ram-fourier}
prove the first assertion.  Lemma~\ref{lem:toral-Fourier},
together with the degree-three torus Stone--von Neumann theorem,
converts those Fourier coefficients into the displayed
multiplicities; their weighted dimensions are $3\cdot9=27$.
For the trivial ordinary character one can choose
$\widetilde\theta$ with $\widetilde\theta(a^*)=1$:
the odd-$n$ scalar star kills $\mu_3\subset F^\times$
(Section~\ref{subsec:star}), and a faithful character of
$\mu_3$ extends across the abelian scalar central group
modulo the star image.  The nontrivial Galois automorphism
acts by inversion on the respective finite quotients' character
parameters (with the $C_2$ factor fixed), so the displayed
multiplicities are invariant.
\end{proof}

\subsection{From toral positivity to the general problem}
\label{subsec:strategywild}
The example suggests two distinct conjectural steps, and shows
why the first alone is insufficient.

\paragraph{Step T: positivity on each field torus.}
For a general prime degree $r$, odd $n$ and wild $F$, test
whether $R_E$ has a nonnegative integral Fourier expansion
on $E^\times/F^\times$ for \emph{every} degree-$r$ field
$E/F$.  A particularly tractable sufficient condition is to
write $R_E$ as a nonnegative integral combination of
\emph{permutation characters} of finite abelian quotients:
\begin{equation}\label{eq:permutation-criterion}
 R_E(t)=\sum_{U\leq A}c_U[A:U]\,\mathbf1_U(t),
 \qquad c_U\in\mathbb Z_{\geq0},
\end{equation}
where $A$ is a finite quotient of $E^\times/F^\times$ and
$U$ ranges over its subgroups.  Indeed
$[A:U]\,\mathbf1_U=\operatorname{Ind}_U^A\mathbf1$ is a
permutation character for abelian $A$.  The two formulas above
are explicit examples of \eqref{eq:permutation-criterion}.
To seek such expressions generally, use the tower of unit
filtrations of $E/F$ and the identity
\begin{equation}\label{eq:cyclo-det}
 R_E(t)=\left|
 \det\!\left(\sum_{j=0}^{n-1}\operatorname{Ad}(t)^j
 \Bigm|\mathfrak d/\mathfrak e\right)
 \right|_F^{-1/2},
\end{equation}
where $\mathfrak d=D$ and $\mathfrak e=E$ as $F$-Lie
algebras; this follows from
$1-X^n=(1-X)(1+X+\cdots+X^{n-1})$
on the noncentral adjoint space.  The finite valuation strata
in \eqref{eq:cyclo-det} are candidates for indicators of
congruence subgroups.  Neither their integrality nor their
positivity is established for general $r,n,F$.

\paragraph{Step G: nonabelian gluing.}
Even if Step~T holds on every maximal field torus, the local
restrictions must arise \emph{simultaneously} from one character
of $\widetilde{D^\times}$.  The ramified example exposes extra
congruence data invisible in the unramified residual quotient:
$R(j)=1$, $R(1+j)=3$, and $R(1+\varpi j)=9$.
Torus positivity does not impose the compatibility of the
corresponding Heisenberg models under conjugation by
$D^\times$ or supply the integral coefficients of the full
character matrix \eqref{eq:matrixchars}.

A concrete test is to choose the maximal order $\mathcal O_D$,
its radical $\mathfrak P_D$, and
$U_D^m=1+\mathfrak P_D^m$; select an $m$ sufficiently large
that the relevant smooth class functions and representations
are controlled by a finite level.  One must retain the genuine
central character in a \emph{twisted} finite group algebra,
rather than quotienting out the genuine center.  At each level:
(1) calculate the induced cocycle and genuine irreducibles;
(2) use Clifford theory for the normal congruence subgroups;
(3) compare their restrictions to all quadratic covered tori;
(4) compute the actual row coefficients
$\langle L_D\mathbf1,\Theta_\tau\rangle$ exactly.
A single noninteger coefficient disproves full integrality;
integer coefficients at one level are not a uniform proof for
all smooth irreducibles.  In the general case one must then
establish level-independent compatibility for arbitrary
ordinary $\rho$, not only $\rho=\mathbf1$.
\section{Connection with Zou's corrected character relation}\label{sec:Zou}
For the Kazhdan--Patterson cover of the \emph{split group},
Zou \cite[\S3]{Zou24} revisits the Flicker--Kazhdan metaplectic
correspondence and records a corrected scalar in the character
relation.  This is an essential input, but its objects are
$\widetilde{\GL_r(F)}$ and ordinary $\GL_r(F)$, not directly
$\tH$ and $H$ at a division place.

\begin{proposition}[\Etag\ Zou's root-sum constant in our case]\label{prop:Zouconstant}
Write $d_F(m)=m|m|_F^{-1/2}$ and
$d_r=\gcd(n,2rc_v+r-1)$.  Under \eqref{eq:parameters},
$d_r=n$.  In Zou's equation (3.2), the coefficient preceding
the root sum in rank $r$ is
\begin{equation}\label{eq:Zouprefactor}
 \frac{1}{d_F(d_r)|n|_F^{r/2}}
 =\frac{1}{n|n|_F^{(r-1)/2}}=\frac{b_D}{n}.
\end{equation}
For regular elliptic $x$ with degree-$r$ field centralizer,
$x^n$ is automatically regular by Remark~\ref{rem:power-singular}.
The equation $t^n=x^n$ then has exactly $n$ solutions
$t=\zeta x$, $\zeta\in\mu_n$.  If the linear character
$\pi$ has central character $\omega_\pi$ trivial on $\mu_n$,
these summands agree and the effective coefficient is $b_D$.
This restriction on $\omega_\pi$ is essential: without it, the
scalar sum is $\sum_{\zeta\in\mu_n}\omega_\pi(\zeta)=0$,
so equation (3.2) instead forces the covered character to vanish
at $x^*$.  For a genuine lift pair the compatibility
$\omega_\pi(a)=\omega_{\widetilde\pi}(a^*)$
and $\zeta^*=1$ imply $\omega_\pi|_{\mu_n}=1$ automatically.
\end{proposition}
\begin{proof}
The equation $d_r=n$ is immediate from scalar centrality;
\eqref{eq:Zouprefactor} follows by substitution.  Regularity of
$x^n$ follows from the field case of Remark~\ref{rem:power-singular}.
Any solution $t$ centralizes $x^n$, so $t\in F[x]^\times$.  In the field $F[x]$, $(t/x)^n=1$,
and $\mu_n(F[x])=\mu_n(F)$.  Thus $t=\zeta x$.
The Weyl discriminant is invariant under central scalar
multiplication, the ordinary character takes the same value
because it is trivial on $\mu_n$, and $(\zeta x)^*=x^*$
by the explicit odd-$n$ scalar star computation in
Section~\ref{subsec:star}.
These are precisely the equal terms in Zou's (3.2).
\end{proof}

\begin{remark}[What is and is not composed]
At the ramified places one may compare the \emph{desired}
$\GDb_v\leftrightarrow\Gb_v$ character relation of Parts I--II
with (a) the classical linear Jacquet--Langlands elliptic
character identity and (b) Zou's formula.  The common coefficient
$b_D$ and the familiar sign $(-1)^{r-1}$ are consistent with
this comparison.  Nevertheless \emph{the existence of an
irreducible genuine image on the covered division side is the
missing conclusion}; combining two known formulas cannot be
used to assume that missing representation exists.
For non-elliptic split centralizers, retain Zou's entire
root sum, including the central-character factors.
\end{remark}

\section{Global geometric comparison: proved input and new obligations}\label{sec:geometric}
There are two comparisons to keep distinct:
\begin{equation}\label{eq:twocomparisons}
 \begin{gathered}
 \GDb(\A_K)\ \longleftrightarrow\ \Gb(\A_K)
 \quad\text{(Parts I--II, under their hypotheses)},\\[3pt]
 D^\times(\A_K)\ \dashrightarrow\ \GDb(\A_K)
 \quad\text{(proposed additional comparison)}.
 \end{gathered}
\end{equation}
The first has a proved simple trace-formula comparison with
regular-support or later matching hypotheses (Part I, Theorem 6.6;
Part II, Theorem 6.2).
The second is not a formal consequence of the first or of the
unitarity of $L_D$.

\subsection{Geometric coefficients and the product formula}
For a regular elliptic rational $\gamma$, with $E=K[\gamma]$,
Part I identifies the torus coefficient $c(\gamma)$, recalled in
\eqref{eq:globalvol}. The same torus measure $dt$ must be used
whenever this term is compared with the corresponding linear
inner-form orbital integral.
Part I uses the normalized orbital integral
$J_v(\gamma,f)=\Delta_v(\gamma)O_v(\gamma,f)$.
For a rational regular element the product formula implies
\begin{equation}\label{eq:discprod}
 \prod_v\Delta_v(\gamma)=1.
\end{equation}
This is because each discriminant is the normalized absolute
value of a nonzero element of $K$ (or its square root), including
the determinant normalization in \eqref{eq:Delta}.
The local \emph{ratio} of discriminants
$\Delta_v(t)/\Delta_v(t^n)$ must be kept until local transfer
and the global product are both defined.

At $v\in S$ Part III, Theorem 9.1, gives the exact local
normalized orbital-integral identity
\begin{equation}\label{eq:localJO}
 J_{\tH_v}\bigl(z\,t^*,\mathcal O_{D,v}f_v\bigr)
 =b_{D,v}\,\widetilde\theta_v(z)^{-1}
 J_{H_v}(t,f_v),
\end{equation}
using the same quotient measures on the two copies of $H_v/Z_v$.
At $v\notin S$, $H_v\cong\GL_r(K_v)$ and the appropriate
root-sum comparison is not \eqref{eq:localJO}; it must use
Flicker--Kazhdan--Zou, and it can involve the kernel
$\mu_n^{k_v}/\mu_n$ of Proposition \ref{prop:power-etale}.
In particular the local factors $b_{D,v}$ for $v\in S$
\emph{cannot} simply be cancelled using
$\prod_v|n|_v=1$ without the split-side contributions.

\subsection{Haar ledger}\label{subsec:ledger}
For $v\in S$, Part II, Definition 2.1, fixes
\begin{equation}\label{eq:formaldegree}
 \vol(D_v^\times/F_v^\times;dg_{D,v})
 =d(\mathrm{St};dg_{\GL,v})^{-1},
\end{equation}
where $d(\mathrm{St};dg_{\GL,v})$ is the formal degree
relative to the chosen $\GL_r(F_v)/F_v^\times$ measure.
The probability measure used in Part III is
$d(\mathrm{St};dg_{\GL,v})\,dg_{D,v}$.
When \emph{both} sides of \eqref{eq:localJO} are scaled by the
same measure, $b_{D,v}$ is unchanged.  Different measures on the
$\GL$ side must still be reconciled with \eqref{eq:formaldegree}.

\begin{center}\small
\begin{longtable}{@{}p{0.27\linewidth}p{0.53\linewidth}p{0.10\linewidth}@{}}
\toprule
Item & Required statement or calculation & Status\\\midrule
\endhead
Ambient cocycle and common center & Part I, Lemma 5.1 and Prop. 5.5; restriction \eqref{eq:before} & [E]\\
Regular rational good locus & Theorem \ref{thm:globalgood} & [P]\\
Global rational/stable classes & Hilbert 90 and Kottwitz's framework; local embeddability into $D$ & [E]\\
Global torus coefficient & Same $dt$ in \eqref{eq:globalvol}, including $|W_\gamma|$ & [E]\\
Local $v\in S$ orbital transfer & Part III, Theorem 9.1, equation \eqref{eq:localJO} & [E]\\
Split-place $v\notin S$ invariance & Lemma \ref{lem:FKnecessary} (necessary); rank-two spherical calculation below & [P]\\
Split-place pointwise transfer and $\lambda_v$ & Proposition \ref{prop:conditionalglobal}; unproved existence & [C]\\
Unramified unit elements & \cite[\S\S11--12]{FK} under $(n,N_r)=1$; transport to Part I cover pending & [E]/[O]\\
Nonregular powers and singular germs & Remark \ref{rem:power-singular}; germ compatibility or support needed & [C]\\
Central identity class & Suppressed in Prop. \ref{prop:conditionalglobal}; separate germ/measure test below & [E]/[C]\\
Spectral expansion & Part I--II packet formula with restrictions; new factor matching & [E]/[C]\\
\bottomrule
\end{longtable}
\end{center}
Here $N_r$ is the least common multiple of the \emph{composite}
positive integers at most $r$ (the convention in [FK]), with $N_r=1$
when there are none.  The classical unramified unit matching in
\cite[\S12, Corollary]{FK} is established under $(n,N_r)=1$
and $p\nmid n$.  Thus $N_2=N_3=1$, and for prime $r=5$
we have $N_5=4$, automatically prime to the odd integer $n$.
Other ranks must retain the published hypothesis.

\subsection{Split-place invariance and a geometric comparison criterion}
At a place $v\notin S$ the algebra $D_v$ is split, but a rational
elliptic field $K[s]$ may decompose locally as a product of fields.
By Proposition~\ref{prop:power-etale}, its local $n$th-root fibre
modulo scalars can contain $n^{k_v-1}$ points.  The global class
$[s]$ supplied by Theorem~\ref{thm:globalgood} is nevertheless
unique.  Therefore a local sum over \emph{all} roots cannot be substituted
for the orbital integral of this one rational representative without
an additional compatibility statement.

\begin{definition}[FK-compatible ordinary functions]\label{def:FKcompatible}
Fix the local central characters $\theta_v,\widetilde\theta_v$
with $\theta_v(a)=\widetilde\theta_v(a^*)$.  A normalized
ordinary test function $f_v$ is called \emph{FK-compatible on the
regular-power locus} if, whenever $x,x'\in H_v$ have $x^n,x'^n$
regular and their star images satisfy
$({x'})^*=z x^*$ for a central lift $z\in\widetilde Z_v$
(after conjugating representatives), one has
\begin{equation}\label{eq:FKinvariance}
 J_{H_v}(x',f_v)
   =\widetilde\theta_v(z)^{-1}J_{H_v}(x,f_v).
\end{equation}
The case $x'=\xi x$ with \emph{scalar} $\xi\in\mu_n(F)$
is automatic, since $\theta_v(\xi)=1$.  The substantive
condition occurs on a product torus
$E_v=\prod_{i=1}^{k_v}E_i$ with $k_v\ge2$:
if $\zeta=(\zeta_i)\in\mu_n(E_v)\simeq\mu_n(F)^{k_v}$
is nonscalar and $x^n$ is regular, then
\begin{equation}\label{eq:nonscalar-root-invariance}
 J_{H_v}(\zeta x,f_v)=J_{H_v}(x,f_v).
\end{equation}
The odd-$n$ star is multiplicative on each lifted torus, and
$\zeta^*=1$, so $(\zeta x)^*=x^*$.  The kernel on the
quotient torus is $\mu_n(F)^{k_v}/\mu_n(F)$, of size
$n^{k_v-1}$ (Proposition~\ref{prop:power-etale}).
Equation \eqref{eq:FKinvariance} is the invariant-function
condition of \cite[\S6]{FK}, expressed with the present
inverse central transformation convention.
\end{definition}

\begin{lemma}[\Ptag\ Necessary split-place invariance]
\label{lem:FKnecessary}
Suppose a genuine $\widetilde f_v$ and a fixed nonzero constant
$\lambda_v$ satisfy
\begin{equation}\label{eq:splitpoint}
 J_{\widetilde H_v}(x^*,\widetilde f_v)
      =\lambda_v J_{H_v}(x,f_v)
\end{equation}
for all $x$ in a conjugation- and scalar-stable region where
$x^n$ is regular, with the identity at central translates understood
with the prescribed inverse-central-character law.  Then $f_v$ obeys
\eqref{eq:FKinvariance} on that region.
\end{lemma}
\begin{proof}
If $({x'})^*=z x^*$, conjugation invariance and the inverse
central transformation law of the genuine test function imply
$J_{\widetilde H_v}(({x'})^*,\widetilde f_v)
=\widetilde\theta_v(z)^{-1}
J_{\widetilde H_v}(x^*,\widetilde f_v)$.
Apply \eqref{eq:splitpoint} twice and divide by the fixed
nonzero $\lambda_v$.  In particular, for $x'=\zeta x$ with nonscalar
$\zeta\in\mu_n(E_v)$, one has $({x'})^*=x^*$ and obtains
\eqref{eq:nonscalar-root-invariance}.  The scalar-root case is
already automatic from the central-character convention.
Thus FK-compatibility is a \emph{necessary} condition,
not an asserted construction of $\widetilde f_v$.
\end{proof}

\begin{proposition}[\Ptag\ Rank-two tame spherical-unit check]
\label{prop:ranktwo-unit}
Let $F$ be nonarchimedean with $\mu_n\subset F$, $p\nmid n$,
$r=2$, and $K_0=\GL_2(\mathcal O_F)$.  Take a spherical ordinary
unit $f^0$, supported on $F^\times K_0$, with an unramified
central character, and the standard quotient measures. Then
$f^0$ obeys \eqref{eq:nonscalar-root-invariance} for all regular
split-torus $x$ with $x^n$ regular.  For a quadratic field torus
$E^\times$, the root-fibre condition modulo scalars is vacuous.
Moreover, when $f^0$ and the genuine spherical unit
$\widetilde f^0$ are transported from the standard
Flicker--Kazhdan cocycle and measure conventions,
\cite[\S12, Corollary]{FK} establishes their regular-power
matching with $\lambda_v=1$ in that normalization.  Whether the
same equality holds \emph{literally} for the selected ambient
cover and global measures of Parts I--IV requires their explicit
splitting and Haar identifications.
\end{proposition}
\begin{proof}
A split regular element is conjugate to
$\operatorname{diag}(a,b)$; multiplying by a scalar, first
consider $a,b\in\mathcal O_F^\times$.  The Iwasawa and
unipotent-orbit calculation, with $\operatorname{vol}(\mathcal
O_F)=1$, yields
\[
 O_{\GL_2}\!\left(\operatorname{diag}(a,b),1_{K_0}\right)
       =|a-b|_F^{-1},\qquad
 \Delta\!\left(\operatorname{diag}(a,b)\right)=|a-b|_F,
\]
so $J=1$.  Multiplication of either eigenvalue by a root of unity
preserves its valuation, hence the normalized orbital integral.
If the two eigenvalue valuations differ, neither $x$ nor such a
translate is conjugate into $F^\times K_0$; both orbital
integrals vanish. Scalar translations contribute only the
specified unramified central-character factor.  This proves the
ordinary split-torus invariance. For a quadratic field $E/F$,
Proposition~\ref{prop:power-etale} gives
$\ker(P_n:E^\times/F^\times\to E^\times/F^\times)=1$.
The final statement invokes the \emph{published} unit matching,
not the elementary split-torus calculation.  Its passage to the
specified ambient cover is a distinct normalization obligation.
\end{proof}

\begin{remark}[Scope of the published fundamental lemma]
\label{rem:fk-range}
The Flicker--Kazhdan result is a theorem for the matching
unramified Hecke units, not an existence result for arbitrarily
prescribed functions at finitely many exceptional places.
Its original hypotheses include $p\nmid n$ and
$(n,N_r)=1$. In rank two the only maximal-torus types are
split and quadratic field tori, so the ordinary invariance
calculation above is exhaustive.  At higher ranks, mixed tori
do not automatically invalidate spherical matching: the
Flicker--Kazhdan theorem covers them as well whenever its
hypotheses apply.
\end{remark}

\paragraph{A regularity and germ obligation.}
Equation~\eqref{eq:splitpoint} is imposed only where $x^n$ is
regular.  On a split torus, $x$ can be regular while $x^n$ is
singular (Remark~\ref{rem:power-singular}); the power map is a
local analytic isomorphism on the torus, but normalized orbital
integrals of a function at images approaching a singular element
may have nonconstant Shalika-germ behavior.  A full transfer
construction must either arrange support away from these
root-collision loci or prove the matching germ conditions there.
The same issue occurs at the complex places, although the cover
there splits: the ordinary power map can still carry regular
elements to singular ones. No universal vanishing of the germ
coefficients is claimed here.

\begin{proposition}[\Ptag\ Conditional geometric comparison with
an explicit normalization factor]\label{prop:conditionalglobal}
Fix the rational splittings and compatible central characters
of Parts I--III, and fix \emph{corresponding} local Haar measures,
global torus measures, and normalizer conventions on the ordinary
and covered sides.  Choose factorizable ordinary and genuine test
functions in a trace-class, regular elliptic support range for
which the two simple trace formulas apply. Assume also that
\emph{every nongood rational noncentral class has zero ordinary
geometric contribution}; for example the test functions may be
chosen to suppress the finitely many such classes in a specified
simple trace-formula support range. Assume:
\begin{enumerate}[label=(\roman*)]
\item For each $v\notin S$, $f_v$ is FK-compatible as in
Definition~\ref{def:FKcompatible}, and a genuine test function
$\widetilde f_v$ satisfies \eqref{eq:splitpoint} for every regular
rational local component $s_v$ of a globally good class $[s^n]$,
with a nonzero constant $\lambda_v$ independent of $s_v$ and of
its local torus type.  The prescribed central-translate identity
holds, $\lambda_v=1$ for almost all $v$, and the hypotheses include
the complex places.
\item For each $v\in S$, use the Part III, Theorem~9.1,
orbital-integral transfer $\widetilde f_v=\mathcal O_{D,v}f_v$,
with coefficient $b_{D,v}=|n|_v^{-(r-1)/2}$ and matching quotient
measures on $H_v/Z_v$ and $\widetilde H_v/\widetilde Z_v$.
The functions obey the reciprocal central transformation laws.
\item The common rational splitting and all local lifts satisfy
the star compatibility $\iota(s^n)=\iota(s)^n$; the global torus
coefficients (including $|W_s|$), central phases, and quotient
measures agree under the stipulated identifications.  The support
condition removes the central identity contribution.  The
ordinary and covered simple geometric expansions are absolutely
convergent in the selected test-function range.
\end{enumerate}
Put
\begin{equation}\label{eq:lambda-product}
 \boxed{\quad
 \Lambda=\Bigl(\prod_{v\notin S}\lambda_v\Bigr)
          \Bigl(\prod_{v\in S}b_{D,v}\Bigr).\quad}
\end{equation}
Then the regular elliptic geometric distributions, and hence the
corresponding simple spectral trace distributions, satisfy
\begin{equation}\label{eq:globalLambda}
 \boxed{\quad I_{\widetilde H}(\widetilde f)
              =\Lambda\,I_H(f).\quad}
\end{equation}
In particular the distributions agree if $\Lambda=1$.
\end{proposition}
\begin{proof}
Since the rational splitting $\iota$ is a homomorphism,
$\iota(s^n)=\iota(s)^n$; at every completion, $\iota(s)_v$
is a lift of $s_v$, so this equals its star lift $s_v^*$.
By Theorem~\ref{thm:globalgood}, globally good noncentral
rational classes are parametrized by the unique global root
class $[s]$ of $[s^n]$, preserving normalizer orders.
The explicit nongood-support hypothesis removes the ordinary
classes without a covered counterpart; the cover has zero
regular orbital integrals on nongood lifts.  Local compatibility
handles the multiple \emph{local} roots without introducing
additional rational summands.  Hypothesis (i) supplies
$\lambda_v J_{H_v}(s_v,f_v)$ off $S$, and (ii) supplies
$b_{D,v}J_{H_v}(s_v,f_v)$ on $S$. By (iii) the remaining torus
coefficients and central phases agree, and the product of Weyl
discriminants is $1$ by \eqref{eq:discprod}.  Every regular
rational term on the covered side is thus $\Lambda$ times the
corresponding ordinary term.  The central identity has been
suppressed, and the stipulated trace-class and convergence
conditions permit summation and the two trace formulas. This
gives \eqref{eq:globalLambda}.  The proof supplies \emph{no}
existence theorem for the matching functions, their exceptional
local constants, or their necessary nongood support restriction.
\end{proof}

\begin{remark}[Normalization versus canonical matching]
The condition $\Lambda=1$ is a normalization once nonzero
matching functions are known: rescaling one $\widetilde f_v$
rescales its $\lambda_v$ and hence $\Lambda$. It is \emph{not}
a theorem about canonically normalized units; rescaling a unit
changes its canonical value.  The meaningful unresolved question
is to compute the constants for the prescribed local transfers,
especially the finitely many exceptional split places, and to
verify the global product in those fixed conventions.
Likewise, compactness of $D^\times(K)Z(\mathbb A_K)\backslash
D^\times(\mathbb A_K)$ eliminates continuous spectrum but does
not by itself justify every unrestricted trace sum; the support
and trace-class conditions above are retained for this reason.
\end{remark}

\subsection{The identity class and an explicit normalization test}
\label{subsec:identity-test}
Because $r$ is prime and $D$ is division, every noncentral
rational element generates a degree-$r$ field and is regular
elliptic.  Thus, after quotienting by the rational center, the
central identity class is the only remaining singular rational
class.  Proposition~\ref{prop:conditionalglobal} suppresses it,
but it can be used as an additional test \emph{if} the test range
is enlarged and the necessary local germ identities are supplied.
\begin{remark}[Central discriminants versus orbital-integral germs]
\label{rem:identity-germs}
Let $d=r(r-1)/2$.  For any torus, as a regular $x\to1$
with $x^n$ regular, the \emph{discriminant ratio} satisfies
\[
 \frac{\Delta(x)}{\Delta(x^n)}\longrightarrow |n|_v^{-d},
\]
because $(1-u^n)/(1-u)\to n$ for each root value $u\to1$.
However this identity \emph{does not by itself yield}
$\widetilde f_v(1)=\lambda_v|n|_v^{-d}f_v(1)$
from \eqref{eq:splitpoint}.  On the split group, normalized
orbital integrals generally have nontrivial Shalika germs near
singular central elements; their limiting coefficients need not
be the evaluation of the test function at $1$.  The requisite
relation between \emph{central test-function values} is therefore
a separate local matching obligation, even when the regular
orbital integrals match.  Compare the unit-specific calculation
of Proposition~\ref{prop:ranktwo-unit} with the general
germ warning following Lemma~\ref{lem:FKnecessary}.
\end{remark}

At $v\in S$, Part III, Corollary~9.3 gives
$\widetilde f_v(1)=c_{D,v}f_v(1)$ with
$c_{D,v}=b_{D,v}|n|_v^{-r(r-1)/2}$.  If, as an \emph{additional unproved local identity-value
hypothesis}, every split place (including complex places) satisfies
$\widetilde f_v(1)=\lambda_v|n|_v^{-d}f_v(1)$ with
$d=r(r-1)/2$, and all denominators are nonzero, the ratio of
global identity-function values becomes
\begin{equation}\label{eq:identityproduct}
 \prod_v\frac{\widetilde f_v(1)}{f_v(1)}
   =\Lambda\prod_v|n|_v^{-r(r-1)/2}
   =\Lambda
\end{equation}
by the product formula.  This does not imply equality of the
\emph{identity trace-formula terms} without a separate choice of
matching quotient measures: homeomorphic adelic quotients need
not have equal volumes with arbitrarily scaled Haar measures.
Nor does it justify retaining the identity term in the simple
trace formula without checking its other hypotheses.

\begin{example}[A possible global numerical ledger]
\label{ex:ranktwoledger}
Let $K=\mathbb Q(\zeta_3)$, $r=2$, $n=3$, and choose a
quaternion division $K$-algebra ramified at the place over $3$
and at one other finite place.  At the wild division place,
$b_{D,v}=3$, $c_{D,v}=27$.  If the finite split places have
matching $\lambda_v=1$, the equation $\Lambda=1$ would require
$\lambda_\infty=1/3$ at the unique complex place.
Since the normalized complex absolute value has
$|3|_\infty=9$, the hypothesized central-value formula would then
give the complex identity factor $(1/3)/9=1/27$, balancing
$27$.  This is a \emph{necessary numerical consistency check}
under the stated assumptions, not a proof of matching complex
test functions or a global comparison.
\end{example}

\subsection{Why Li, Kottwitz, and DKV do not close the gap}\label{subsec:litroute}
The linear Jacquet--Langlands theorem of
Deligne--Kazhdan--Vign\'eras \cite{DKV} controls the ordinary
$D^\times$--$\GL_r$ relation, including the elliptic character
sign, but it does not automatically manufacture a genuine
representation of $\GDb_v$.  Kottwitz's analysis of stable
conjugacy and relative trace coefficients \cite{Kottwitz84}
provides the correct arithmetic setting; for the induced tori
$\operatorname{Res}_{E/K}\mathbb G_m$ one has
$H^1(K,T)=1$ by Shapiro's lemma and Hilbert 90, and embeddings
$E\hookrightarrow D$ are controlled by Brauer splitting.
Wen--Wei Li proves an invariant trace formula for adelic covers
under an invariant trace Paley--Wiener hypothesis for Levi subgroups
at \emph{real} archimedean places \cite{Li14}.  Here $K$ is totally
imaginary, so that real-place hypothesis is vacuous.  At the complex
places the specified covers split (Part I, Appendix~B, item~(R4);
see also its ambient-cover construction), and the linear invariant Paley--Wiener theorem
of Clozel--Delorme \cite{ClozelDelorme} supplies the relevant
archimedean background.  Li also discusses known real metaplectic
$\GL$ cases, but his general result does not assert the particular
\emph{equality} of ordinary-$D^\times$ and covered-$D^\times$
traces required in \eqref{eq:twocomparisons}.  The simpler already-proved trace formulas
of Parts I--II remain the appropriate initial comparison.

\section{Spectral separation and a research program}\label{sec:spectral}
\subsection{The actual spectral distribution identity}
Let $H=D^\times$ and $\widetilde H=\GDb$ globally, with
compatible central characters $\theta,\widetilde\theta$.
The quotient $H(K)Z_H(\A_K)\backslash H(\A_K)$ is compact
because $D$ is a division algebra over $K$; its covered analogue
is compact modulo the specified central preimage as well.
Thus the fixed-central-character right regular representations
have discrete Hilbert decompositions.  For factorizable test
functions with the inverse central transformation law, define
\begin{align}
 I_H(f)&=\sum_{\rho\in\mathcal A(H,\theta)}
       m_H(\rho)\prod_v\operatorname{tr}\rho_v(f_v),\\
 I_{\widetilde H}(\widetilde f)
     &=\sum_{\widetilde\rho\in\mathcal A(\widetilde H,\widetilde\theta)}
       m_{\widetilde H}(\widetilde\rho)
       \prod_v\operatorname{tr}\widetilde\rho_v(\widetilde f_v).
       \label{eq:spectralglobal}
\end{align}
Here $\mathcal A$ denotes the relevant automorphic spectrum
and $m$ is its actual spectral multiplicity, not an assumed $1$.
The expressions are interpreted on trace-class test functions
satisfying the simple trace-formula hypotheses; a bare
unrestricted infinite sum is not justified by notation alone.
Under the assumptions of Proposition~\ref{prop:conditionalglobal},
the comparison of geometric expansions supplies
$I_{\widetilde H}(\widetilde f)=\Lambda I_H(f)$;
ordinary equality holds when the fixed normalizations give
$\Lambda=1$.  It does not
identify summands.  This is a conditional identity of
\emph{distributions}, with both the multiplicities and local
trace factors retained.

\subsection{Individual local characters and integrality}
An equality of global spectral sums is a statement about
\emph{distributions}.  To isolate a single local irreducible at a
wild place, one needs an argument that the relevant coefficients
are integers (or multiplicities which can be shown to isolate a
single constituent).  Linear independence alone can separate
finite lists of characters; it does not supply the required
integrality from a unitary Fourier transform.

\begin{remark}[What an integer spectral coefficient would prove]
If a justified global-to-local argument supplies
$a_{\rho\tau}\in\mathbb Z$ for every $\tau$ in the row of a fixed
$\rho$, Proposition~\ref{prop:int} (Part III, Corollary~5.2)
forces $L_{D,v}\Theta_\rho$ to be one irreducible genuine
character with positive sign.  If integrality holds for all
rows, unitarity implies a bijection.  This is an application of
an already proved lattice criterion, not a second conditional
theorem; the missing step is precisely the integrality
hypothesis at a wild place.
\end{remark}

\paragraph{Separating the remaining tasks.}
First, the toral Fourier results in
Proposition~\ref{prop:all-three-tori} establish a necessary
integrality property on every quadratic maximal torus in the
rank-two wild example, but do not solve the nonabelian gluing
problem of Subsection~\ref{subsec:strategywild}.  The latter
requires an actual twisted finite character table and a proof
that the chosen finite level captures the desired row.  Second, at split
places the root sum in Zou's (3.2) should be written against the
kernel and index in Proposition \ref{prop:power-etale}, including
the chosen Haar measures and central phases.  Third, one needs the explicit ambient-cover identification
that transports the known unramified unit matching of
\cite[\S12]{FK} to the selected convention, together with
functions at the finitely many exceptional split places and a
complete global coefficient identity as in
Proposition~\ref{prop:conditionalglobal}. The ordinary test
functions must also suppress nongood rational classes in the
chosen global support range.
Finally, one must justify globalization and spectral separation
for representations of the covered inner form, and remove the
auxiliary supercuspidal and non-cancellation hypotheses of Parts
I--II.  These are mathematical work packages, not restatements of
an already-proved theorem.

\section{Selected glossary}\label{sec:glossary}
\begin{description}[style=nextline,leftmargin=3.0cm]
\item[Good element] A regular lift commuting with the full preimage of its ordinary centralizer; genuine invariant characters vanish on the nongood regular locus.
\item[Scalar-central cover] A cover whose entire inverse image of $F^\times I_r$ is central; here forced by $n\mid r-1+2rc_v$.
\item[Power-map kernel] $\ker(P_n)$ on $T(F)/F^\times$, distinct from the image index.  The kernel can be nontrivial on product tori even when the corresponding field-torus map is injective.
\item[Star map] $h^*=\widetilde h^{\,n}$, independent of lift; multiplication on tori uses $n$ odd.
\item[Weyl discriminant] The factor $\Delta=|\det(1-\operatorname{Ad}(t)|\mathfrak g/\mathfrak t)|^{1/2}$; it is not the unnormalized characteristic-polynomial discriminant.
\item[Toral Fourier test] Expand the discriminant ratio $R_E$ in characters of the compact abelian group $E^\times/F^\times$.  Integral positive coefficients ensure torus-restriction positivity, not full-group integrality.
\item[Nonabelian gluing] The additional requirement that compatible genuine covered-torus characters arise from an integral character of the full group $\widetilde{D^\times}$.
\item[Twisted finite quotient] A finite quotient of $H/Z$ with the central character retained as a multiplier, used to compute projective/genuine irreducible characters.
\item[Packet identity] An identity of traces over collections of automorphic representations; it does not assert a bijection between individual local irreducibles.
\item[Formal degree] The Plancherel normalization of a square-integrable representation relative to a specified quotient Haar measure.
\item[Root-sum relation] Zou's equation (3.2), which sums over all $n$th roots satisfying the stated regularity condition rather than choosing one arbitrarily.
\end{description}

\section*{Acknowledgments}
This Part~IV manuscript was prepared with mathematical and editorial
assistance from ChatGPT (OpenAI).  Claude (Anthropic) supplied
referee reports on Drafts~1--3, including the correction to the
Haar-modulus proof and the split-place matching obligations.  Earlier Parts I--III were developed
with separate assistance acknowledged there.  The author contributions and the referee corrections are documented
in the separate revision log.  The rank-two Fourier calculations and
conditional proof strategy were added in Draft~5 after a separate
comparison of Parts~I--III and [JC].  The remaining research problems are
recorded in the status table and Section~\ref{sec:spectral}.

\begin{thebibliography}{99}
\bibitem{DKV} P. Deligne, D. Kazhdan and M.-F. Vign\'eras,
\emph{Repr\'esentations des alg\`ebres centrales simples $p$-adiques},
in \emph{Repr\'esentations des groupes r\'eductifs sur un corps local},
Travaux en Cours 8, Hermann, Paris (1984), 33--117.
\bibitem{FK} Y. Flicker and D. Kazhdan,
\emph{Metaplectic correspondence}, Publ. Math. IH\'ES \textbf{64}
(1986), 53--110, especially \S\S2--4, 12, 18, 24--28.
\url{https://www.numdam.org/articles/10.1007/BF02699192/}.
\bibitem{Zou24} J. Zou,
\emph{Local metaplectic correspondence and applications},
arXiv:2209.05852v3 (19 December 2024), especially
(2.2), (3.2), Remarks 3.1 and 3.8, and Proposition 3.6.
\url{https://arxiv.org/abs/2209.05852}.
\bibitem{KP} D. Kazhdan and S. Patterson,
\emph{Metaplectic forms}, Publ. Math. IH\'ES \textbf{59}
(1984), 35--142.
\bibitem{Kottwitz84} R. E. Kottwitz,
\emph{Stable trace formula: cuspidal tempered terms}, Duke Math. J.
\textbf{51} (1984), 611--650.
\bibitem{ClozelDelorme} L. Clozel and P. Delorme,
\emph{Le th\'eor\`eme de Paley--Wiener invariant pour les groupes de Lie
r\'eductifs}, Invent. Math. \textbf{77} (1984), 427--453.
\bibitem{Li14} W.-W. Li,
\emph{La formule des traces pour les rev\^etements de groupes
r\'eductifs connexes. IV. Distributions invariantes},
Ann. Inst. Fourier \textbf{64} (2014), 2379--2448.
\url{https://doi.org/10.5802/aif.2915}.
\bibitem{JC} D. Joyner with Claude (Anthropic),
\emph{On the metaplectic analog of Kazhdan's ``endoscopic'' lifting},
corrected 2026 version of the 1988 paper,
\texttt{metaplectic\_kazhdan\_revised12.tex}, especially \S\S2, 6.
\bibitem{JoynerI} D. Joyner, Claude (Anthropic), and ChatGPT (OpenAI),
\emph{A metaplectic Jacquet--Langlands transfer for division algebras
of prime degree, I: Local matching and a global transfer under
elliptic-support conditions}, Draft 8 (2026), especially \S\S2--6.
\bibitem{JoynerII} D. Joyner, Claude (Anthropic), and ChatGPT (OpenAI),
\emph{A metaplectic Jacquet--Langlands transfer for division algebras
of prime degree, II: Matching of arbitrary functions and an improved
global transfer}, Draft 8 (2026), especially \S\S2, 5--7.
\bibitem{JoynerIII} D. Joyner, ChatGPT (OpenAI), and Claude (Anthropic),
\emph{A metaplectic Jacquet--Langlands transfer for division algebras
of prime degree, III: The division-algebra side and its local
character transform}, Draft 9 (2026), especially Corollary~7.8
and \S\S3--5, 7--9, 11--12.
\end{thebibliography}
\end{document}
